\chapter{界面、导航与回放}
\section{表现层组织}
\code{main.tscn}加载以\code{game.gd}为核心的界面。大量页面使用\code{_draw()}与自定义热点；输入、弹窗、阅读器和稳定列表采用原生控件。建设、科研、休整分别由\code{construction_list.gd}、\code{research_list.gd}和\code{rest_summary.gd}辅助。

资源、工业、指挥科研区没有各自独立的Godot场景文件，而是同一控制器中的页面状态、背景和热点布局。产品上的“进入场景”，在代码层面是页面状态切换。

\section{导航闭环与草稿}
导航历史保存页面、建筑、车间、模式和基地右栏，最多32项；返回恢复区域上下文。科研详情返回原清单时控件继续存在并保留位置。主线与核心采用相同位置切换，各自章节选择独立保存。

\code{draft}与持久\code{formation}分离。预设和出战按可用库存重新分配。部署保存原草稿，确认时将编队与出击放入同一命令事务，失败不会仅覆盖编队；待响应时阻止重复提交。

\section{紧凑清单与资源顶栏}
两种\code{ScrollContainer}行按ID复用；revision、请求状态或字号变更时更新，倒计时不重建整列。行保留标题、状态、收益/工时、两行材料、操作和详情。活动项同位置切换金币加速，科研等待项不可越过前项加速。

顶栏显示库存、保护线、容量和产率，达到容量后自产暂停。K/M/G/T只用于显示，计算和输入保持精确整数。

\section{动画与事件流}
\code{actions}表示完整攻击，\code{events}表示逐发命中、暴击、未命中或空地。Godot按动作播放，支持暂停、倍速、步进、跳过与重播，使用事件余量而不重算伤害。

坦克与歼击车平射，自行火炮与火箭使用发射段、高角度下落段。火箭空位仍落地留坑，地面痕迹与车辆损伤分层。发射声按整次齐射播放，命中声采用紧凑序列；声音池减少发射、爆炸和击毁互相覆盖。

\impl{弹道、声音和换手停顿属于播放器。若新增飞行时间影响命中等玩法，需要另改规则与事件协议，不能只改动画。}
\src{native/scripts/game.gd}{_navigate；_battle_step；_next_fire_interval；open_report}
